\documentclass{article}
\usepackage[preprint]{neurips_2026}
\usepackage[utf8]{inputenc}
\usepackage[T1]{fontenc}
\usepackage{hyperref}
\usepackage{url}
\usepackage{booktabs}
\usepackage{amsfonts,amsmath,amssymb,amsthm,mathtools}
\usepackage{microtype}
\usepackage{xcolor}
\usepackage{graphicx}
\usepackage{enumitem}

\newcommand{\etaL}{\ensuremath{\eta/L}}
\newcommand{\rhoG}{\ensuremath{\rho_G}}
\newcommand{\R}{\mathbb{R}}
\newtheorem{proposition}{Proposition}

\title{Invariance is not robustness: a margin-to-Lipschitz view of foundation models}

\author{Anass Al Ammiri\\
  Technical University of Munich\\
  \texttt{anass.al-ammiri@tum.de}}

\begin{document}
\maketitle

\begin{abstract}
A common intuition holds that a model which is more invariant to nuisance changes of its input, such
as small shifts of an image or rewordings of a prompt, is also more robust to adversarial attacks. We
show that this intuition does not survive contact with strong attacks in foundation models. Working
with a threat-matched margin-to-Lipschitz ratio \etaL{}, defined as the mean classification margin on
clean inputs divided by the mean dual-norm of the input gradient of that margin, we find a clean
dissociation. Through frozen vision-language encoders, \etaL{} computed on clean images alone, without
running any attack, ranks eleven towers by their AutoAttack robustness and orders the per-image
certified radius, while shift-consistency does neither. The ranking survives removing the effect of
adversarial training as a covariate, whereas a cosine-similarity form of consistency does not: it only
appeared to work because it detects adversarial training. In instruction-tuned language models the two
notions become measurably independent. A sensitivity axis (the same ratio \etaL{}) and an
excessive-invariance axis (an orbit-flip radius \rhoG{}, the size of the smallest meaning-changing edit
a model wrongly ignores) are uncorrelated across prompts, and we relate \rhoG{} to a certificate on the
model's robust radius. We also show that the apparent link between invariance and robustness is an
artifact of weak attacks: on every non-robust vision-language tower, single-step FGSM reports
16--46\% robust accuracy that AutoAttack erases to zero. We close by reporting, honestly, where the
transfer stops.
\end{abstract}

\section{Introduction}
\label{sec:intro}
% CONTRIBUTIONS (locked):
% 1. VLM: eta/L is an attack-free cross-tower AND per-image robustness predictor; shift-consistency is not.
% 2. LLM: sensitivity (eta/L) and invariance (rho_G) are orthogonal; rho_G certifies a robust radius.
% 3. weak-attack artifact in frozen VLM towers (FGSM overstates, AutoAttack erases).
% 4. honest scope: encoder-level; first-token margin does not predict multi-token jailbreak radius (matched-margin test).
% Define eta/L, threat-matched, gauge-free, shift-consistency, rho_G here at first use.
A model that assigns the same output to a shifted image, or the same answer to a reworded prompt, is
often described as robust, and this invariance is frequently offered as an explanation for, or a route
to, adversarial robustness. The two properties are not the same. Invariance is about ignoring changes
that should not matter; adversarial robustness is about not being flipped by a small worst-case change.
This paper asks whether invariance predicts adversarial robustness in modern foundation models, and
finds that it does not, while a simple margin-based quantity does.

That quantity is a threat-matched margin-to-Lipschitz ratio. For an input $x$ with correct label $y$,
let the margin be $M(x) = f_y(x) - \max_{j\neq y} f_j(x)$, the gap between the correct logit and the
strongest competitor. Let $\eta$ be the mean of $M$ over correctly-classified clean inputs, and let $L$
be the mean of $\lVert \nabla_x M \rVert_q$, the size of the input gradient of that margin measured in
the norm dual to the attack ($q=1$ for an $\ell_\infty$ attack, $q=2$ for an $\ell_2$ attack). The
ratio $\etaL$ has two properties that make it the right scalar. It is threat-matched, since the dual
norm is chosen to match the attack. And it is gauge-free: multiplying all logits by a constant $c$
scales $\eta$ and $L$ by the same $c$, so the ratio is unchanged, which means a model's logit
temperature cannot inflate or deflate it. A first-order argument gives the certificate that the robust
radius of $x$ is at least $\etaL$, so a larger ratio should mean a more robust input.

We make four claims. First, computed through frozen vision-language towers on clean images alone,
$\etaL$ ranks eleven towers by their AutoAttack robustness and orders the per-image certified radius,
while shift-consistency does neither; the ranking survives removing adversarial training as a
covariate, and a cosine form of consistency does not. Second, in instruction-tuned language models the
sensitivity axis $\etaL$ and an excessive-invariance axis, an orbit-flip radius $\rho_G$, are
measurably independent, and $\rho_G$ upper-bounds the model's oracle-robust radius. Third, the apparent
link between invariance and robustness is an artifact of weak attacks: on every non-robust tower,
single-step FGSM reports robust accuracy that AutoAttack erases, a gradient-masking signature. Fourth,
we report honestly where the transfer stops, at the boundary between an encoder-level margin and a
multi-token generation attack.

\section{Related work}
\label{sec:related}
% Must-cite + one-line distinction each (see SYNTHESIS_lit_and_framing.md, verified with verbatim quotes):
% Attack-free ranking: GREAT Score (2304.09875, generative-manifold, plain classifiers), RDI (2504.18556, feature-cluster).
% Margin/Lipschitz predictors: Ngnawe (2406.18451, within-model per-sample bare margin, no L); MaCS (2603.05812, training regularizer, noise/blur consistency not shift, no towers).
% Invariance vs robustness: Singla-Ge (2103.02695, causal margin, vision, no eta/L, no AT, no strong-attack artifact); Tramer 2020 (2002.04599, population tradeoff, not measured orthogonality).
% Certified NLP (opposite direction): CERT-ED 2408.00728, RS-Del 2302.01757, Text-CRS 2307.16630.
% Weak-attack / gradient masking: Athalye 2018 (1802.00420) parent phenomenon.
% VLM robustness governed by tower: 2407.11121; FARE 2402.12336; anti-aliasing Zhang 2019 (1904.11486).
\textbf{[DRAFT PENDING.]}

\section{Background: the dissociation on image classifiers}
\label{sec:cifar}
% One paragraph + cite prior CIFAR work: eta/L predicts robustness (Pearson ~0.998 across arms),
% shift-consistency anti-predicts under AT (Linf -0.88), threat-matched eta/L predicts AutoAttack (+0.90);
% FGSM shows invariance helping, AutoAttack erases it. Sets up the generalization to foundation models.
The dissociation was first established on image classifiers, which we recap briefly as the origin of
the present study. Across a family of CIFAR-10 models that vary only in how much shift-invariance they
build in, the threat-matched $\etaL$ tracks robustness closely (Pearson around $+0.998$ under an
$\ell_2$ threat model on standard-trained arms), while shift-consistency is uninformative to
anti-predictive. Under adversarial training the split is sharper: shift-consistency anti-predicts
AutoAttack robustness (around $-0.88$ under an $\ell_\infty$ threat model), and the exactly invariant
arm, which has perfect consistency, is the least robust, while the threat-matched $\etaL$ still
predicts (around $+0.90$). The same two patterns hold across MNIST, Fashion-MNIST, a ResNet backbone,
and ImageNet-100, where clean accuracy is flat across arms so the difference is margin geometry rather
than underfitting. The natural question, which this paper answers, is whether these patterns are a
property of small convolutional classifiers or of learned representations more broadly.

\section{Vision-language encoders}
\label{sec:vlm}
We first ask whether \etaL{} predicts adversarial robustness across frozen vision-language towers,
and whether shift-consistency does. This is the setting where the two notions are most often
conflated, because an encoder that assigns nearby representations to shifted views of an image is
routinely described as robust.

\paragraph{Setup.}
We take a panel of eleven frozen vision towers used for zero-shot classification on ImageNet-100:
four adversarially trained encoders (FARE at two radii, TeCoA at two radii), a base CLIP tower, a
DINOv2 tower, and five further CLIP variants spanning pretraining corpus, patch size, and capacity.
For each tower we compute, on clean correctly-classified images only, the margin $M(x)$ of the
zero-shot logits, the ratio \etaL{} with $L$ the mean $\ell_1$ norm of $\nabla_x M$ (the dual norm of
an $\ell_\infty$ attack), and a shift-consistency score, the fraction of integer pixel shifts under
which the top-1 label is unchanged. We then measure adversarial robustness directly with AutoAttack,
reported as the APGD-CE plus APGD-T ensemble, both as a per-tower robust accuracy and as a per-image
certified radius from a fine grid. All quantities use the same set of clean, correctly-classified
images per tower, and \etaL{} is gauge-free, so a tower's logit temperature cannot affect it.

\paragraph{Across towers, \etaL{} predicts robustness and consistency does not.}
The ratio \etaL{} correlates with per-tower robust accuracy at Pearson $+0.953$
($95\%$ CI $[0.92, 1.0]$, permutation $p=0.0009$). The concern with an eleven-tower panel is that any
quantity separating adversarially trained from ordinary towers will correlate with robustness, since
the four trained towers carry almost all of the robustness. We therefore remove an adversarial-training
indicator as a covariate. The partial correlation of \etaL{} with robustness, controlling for that
indicator, is $+0.51$, and controlling for both the indicator and clean accuracy it is $+0.56$: the
ratio carries robustness information that ``this tower was adversarially trained'' does not. A
cosine-similarity form of consistency looks strong at first ($+0.81$) but collapses to $+0.15$ under
the same control, because it works only by detecting adversarial training. The prediction-agreement
form of shift-consistency, the version the anti-aliasing literature tracks, does not order robustness
at all (Pearson $+0.37$, not significant). Used as a selection rule, picking the tower with the
largest \etaL{} costs 5 points of robust accuracy against the oracle, while picking by
prediction-agreement consistency or by clean accuracy costs 88 points.
% Fig: results/c1_tower/figures c1_dissociation_main + a partial-correlation bar (regenerate for n=11).

\paragraph{Per image, the dissociation is sharper and is not an artifact of adversarial training.}
The powered version of the claim is per image. Pooling correctly-classified images across the four
adversarially trained towers, the per-image \etaL{} orders the per-image robust radius at Spearman
$+0.78$ ($95\%$ CI $[0.75, 0.81]$, $n=1200$), while per-image shift-consistency orders it at
$+0.12$ ($[0.06, 0.18]$). With a finer radius grid the ordinary CLIP towers also acquire a real spread
of per-image radii, and within each of them \etaL{} orders the radius at Spearman between $+0.54$ and
$+0.79$, every confidence interval above zero. The per-image dissociation therefore holds inside
ordinary, non-robust towers as well, so it is not a byproduct of the trained-versus-untrained split.
% Fig: per_image_*.pt -> the per-image eta/L vs radius scatter (AT pooled + within non-AT). Headline figure.

\paragraph{What survives and what to state honestly.}
The implementation was checked line by line and no result changes sign under any check; there is no
gradient masking, since APGD $\leq$ PGD-40 $\leq$ FGSM and a black-box Square attack does not beat
APGD on any tower. The honest limitation is that the seven non-robust towers all sit at essentially
zero robust accuracy under strong attack, so the range of tower-level robustness comes from the
trained towers; the covariate control removes the confound in the ranking, and the powered claim rests
on the per-image axis, which stands on both trained and ordinary towers.

\section{The invariance-robustness correlation is a weak-attack artifact}
\label{sec:weakattack}
A recent line of work argues that architectural invariance or equivariance confers adversarial
robustness on its own, without adversarial training, and reports this using single-step FGSM and
short PGD. One representative paper attributes the effect to gradient regularization and smoother
decision boundaries and never runs a strong attack. We show that on frozen, non-adversarially-trained
vision-language towers this apparent link is an artifact of weak attacks.

\paragraph{Setup and the diagnostic we apply.}
The standard way to catch a false sense of robustness is the gradient-masking checklist: a single-step
attack should not beat an iterative one, robust accuracy should fall as attack iterations increase, and
a black-box attack should not beat a white-box one. We run, on each of the seven non-robust towers and
the four adversarially trained towers as a contrast, FGSM and PGD at 2, 5, 10, 20, and 40 steps, then
the APGD-CE plus APGD-T AutoAttack ensemble and a black-box Square attack, all at $\varepsilon=2/255$.
For each attack strength we also correlate the tower's shift-consistency with its robust accuracy
across towers, which is the quantity the pro-invariance claim relies on.

\paragraph{Result.}
On every non-robust tower, single-step FGSM reports $16$ to $46\%$ robust accuracy, which decays as
PGD iterations increase and reaches $0.0$ under AutoAttack; the black-box Square attack does not beat
APGD, and the four adversarially trained towers do not collapse. The correlation between
shift-consistency and robust accuracy across towers, which is positive and sizeable when robustness is
measured with FGSM, decays toward zero as the attack strengthens and is not present under AutoAttack.
The ``more invariant is more robust'' ordering therefore exists only while the attack is weak; it is a
gradient-masking signature, not a property of the encoders. We report this as a targeted correction of
the specific pro-invariance claim, not as a discovery that undefended encoders are fragile, which is
already known.
% Fig (run_weakattack.py, pending free GPU): (a) PGD-k monotonicity per tower; (b) corr(consistency, robust acc) vs attack strength -> null at AutoAttack.
% Must-cite: 2510.16171 (target), 1802.00420, 1902.06705, 2003.01690, 2402.12336, 2401.04350, 2103.02695, 2307.09804.

\section{Instruction-tuned language models}
\label{sec:llm}
% LOCKED except final framing. Sources: verify/B2_verification.md §3e.
\paragraph{Sensitivity and excessive invariance are independent axes.}
We move to instruction-tuned language models (Llama-3-8B-Instruct, with Qwen-2.5-7B as a contrast) and
a corpus of 1600 items across sentiment, negation natural-language inference, and harmful-request
tasks, each with a deterministic, non-circular oracle for a meaning-changing edit. Two robustness
notions are available here. The sensitivity ratio \etaL{}, computed with the task-appropriate class
margin, measures how far a correct decision sits from the boundary relative to how fast the margin
moves. The excessive-invariance radius \rhoG{} measures the size of the smallest meaning-changing edit
(an antonym swap, a negation, or a harmful-to-benign rewrite) that the model wrongly answers the same
way. Across items these two are uncorrelated: Spearman$(\etaL, \rhoG) = -0.025$ ($95\%$ CI
$[-0.09, +0.04]$). Sensitivity to small worst-case perturbations and invariance to large
meaning-changing edits are, in these models, orthogonal axes.
% Fig: results eta/L vs rho_G scatter, null.

\paragraph{An orbit-flip radius and its certificate.}
\begin{proposition}[informal]
\label{prop:rhoG}
For a prompt $x$ with oracle label $y(x)$, let \rhoG$(x)$ be the size of the smallest edit that
changes the meaning yet leaves the model's answer unchanged. Then \rhoG$(x)$ is an upper bound on the
oracle-robust radius of $x$: no smaller edit both changes the meaning and is answered correctly.
\end{proposition}
This inverts the direction of the certified-NLP literature, which certifies that a prediction is
\emph{unchanged} within a radius as a virtue; here the same geometry certifies that staying unchanged
past \rhoG{} is a failure. We report the distribution of \rhoG{} in token-edit and embedding distance
rather than the tautological statement that no meaning-changing edit is smaller than the smallest one.

\paragraph{Excessive invariance, decomposed honestly.}
The pooled rate at which the model wrongly ignores a meaning-changing edit is $0.564$, but this pools
two very different phenomena and we do not report it as a single number. On sentiment, where the model
is a genuine classifier (it answers negative on $54\%$ and positive on $46\%$ of inputs, at $92\%$
accuracy), it wrongly ignores a meaning-changing edit on $0.67$ of the items it gets right, and on the
clean antonym subset alone on $0.673$; this is genuine excessive invariance. On negation inference and
on harmful requests the same rate is instead a constant-classifier degeneracy: the model answers
``entailment'' on all $500$ inference items regardless of the gold label, and refuses $97\%$ of
requests, so a high invariance rate simply re-encodes the label prior. This degeneracy is
model-dependent, since Qwen-2.5-7B is not constant on the same corpus. We treat the constant-classifier
case as a boundary condition and relate it to the known observation that a constant classifier attains
an unbounded margin ratio with no real robustness.
% Table: per-family flip decomposition (verify/B2_verification.md §3e).

\section{Where the transfer stops}
\label{sec:tdiss}
% PENDING the matched-margin run (fix_matched_margin.py) + masking/GCG.
% Preliminary: a FIRST-TOKEN refuse margin does NOT predict the multi-token L2 jailbreak radius
% (Spearman -0.17, partial -0.067). Gauge-invariance mechanism confirmed. Diagnosed as a threat-model
% mismatch (first-token margin vs multi-token generation attack); matched-margin R_cont=L0/||grad L0||
% test resolves whether it transfers when properly matched. FINALIZE after that run + masking/GCG.
\textbf{[DRAFT PENDING matched-margin result. Honest reporting: a first-token refuse margin does not
predict the multi-token jailbreak radius; the gauge-invariance mechanism is confirmed; we identify the
threat-model mismatch and report the threat-matched test.]}

\section{Discussion}
\label{sec:discussion}
\textbf{[DRAFT PENDING. Limitations: encoder-level vs decision-boundary transfer; model-dependence of
the degeneracy; per-image is the powered axis on the tower side; embedding geometry of \rhoG{} is
coarse.]}

\bibliographystyle{plainnat}
% \bibliography{refs}   % build refs.bib from SYNTHESIS_lit_and_framing.md must-cite lists
\end{document}
